\section*{第 \ref{ch:g12:respiration-fermentation} 章——细胞呼吸与发酵}

\begin{solution}{exo:g12:respiration-fermentation:1}
一种带有三个相连磷酸基团的\omterm{def:g10:universal-dna:nucleotide}{核苷酸}，是\omterm{def:g10:cells-common-unit:cell}{细胞}的能量
“货币”。使用时去掉末端磷酸，释放约
\qty{30}{kJ} 每摩尔能量，用以驱动某一过程，并留下 ADP
待再次充能。
\end{solution}

\begin{solution}{exo:g12:respiration-fermentation:2}
\omterm{prop:g12:respiration-fermentation:glycolysis}{糖酵解}在\omterm{def:g10:cells-common-unit:cell}{细胞质}中进行：生成 2 个 \omterm{def:g12:respiration-fermentation:atp}{ATP} 和 2 个还原型 NAD，以及两分子丙酮酸。
\omterm{prop:g12:respiration-fermentation:mitochondrion}{克雷布斯循环}在\omterm{def:g10:cells-common-unit:organelle}{线粒体}基质中进行：放出二氧化碳，生成
还原型载体及 2 个 \omterm{def:g12:respiration-fermentation:atp}{ATP}。呼吸链位于\omterm{def:g10:cells-common-unit:organelle}{线粒体}内膜：还原型
载体将电子交给氧，生成水，并产生约 26 个 \omterm{def:g12:respiration-fermentation:atp}{ATP}。
\end{solution}

\begin{solution}{exo:g12:respiration-fermentation:3}
\omterm{def:g10:cells-common-unit:organelle}{线粒体}有两层膜，内膜向内折叠形成嵴，包围
基质。基质中进行\omterm{prop:g12:respiration-fermentation:mitochondrion}{克雷布斯循环}；内膜上是呼吸链
和 \omterm{def:g12:respiration-fermentation:atp}{ATP} 合成。
\end{solution}

\begin{solution}{exo:g12:respiration-fermentation:4}
它把\omterm{prop:g12:respiration-fermentation:glycolysis}{糖酵解}产生的还原型 NAD 所携带的氢交给丙酮酸，使
\omterm{prop:g12:respiration-fermentation:glycolysis}{糖酵解}无需氧即可继续进行。产量仅为\omterm{prop:g12:respiration-fermentation:glycolysis}{糖酵解}的 2 个 \omterm{def:g12:respiration-fermentation:atp}{ATP}，
另加乳酸或乙醇和 $\mathrm{CO_2}$。
\end{solution}

\begin{solution}{exo:g12:respiration-fermentation:5}
呼吸作用约产生 30 个 \omterm{def:g12:respiration-fermentation:atp}{ATP}，\omterm{def:g10:cell-metabolism:fermentation}{发酵}只产生 2 个。呼吸作用的产物是二氧化碳和水；
\omterm{def:g10:cell-metabolism:fermentation}{发酵}的产物是乳酸，或乙醇和二氧化碳，其中仍储存着大量
能量。
\end{solution}

\begin{solution}{exo:g12:respiration-fermentation:6}
葡萄糖：无影响——\omterm{def:g10:cells-common-unit:organelle}{线粒体}不能利用它。丙酮酸：耗氧开始——
它是\omterm{def:g10:cells-common-unit:organelle}{线粒体}的\omterm{def:g11:enzymes-and-phenotype:enzyme}{底物}。ADP：耗氧加快——耗氧与
\omterm{def:g12:respiration-fermentation:atp}{ATP} 合成相偶联。氰化物：耗氧停止——呼吸链
是耗氧的场所。
\end{solution}

\begin{solution}{exo:g12:respiration-fermentation:7}
\omterm{prop:g12:respiration-fermentation:glycolysis}{糖酵解}将葡萄糖转化为丙酮酸，在\omterm{def:g10:cells-common-unit:cell}{细胞质}中进行；
\omterm{def:g10:cells-common-unit:organelle}{线粒体}缺乏相应的酶，它摄入的是丙酮酸，
而非葡萄糖。
\end{solution}

\begin{solution}{exo:g12:respiration-fermentation:8}
每分子葡萄糖产生两个还原型 NAD；\omterm{def:g10:cells-common-unit:cell}{细胞}内的 NAD 储量很少，
一旦全部还原，\omterm{prop:g12:respiration-fermentation:glycolysis}{糖酵解}便没有载体接收
所释放的氢而停止。\omterm{def:g10:cell-metabolism:fermentation}{发酵}将 NAD 上的氢交给丙酮酸，
使这一途径得以持续。
\end{solution}

\begin{solution}{exo:g12:respiration-fermentation:9}
呼吸作用：$60/30 = 2$ 分子葡萄糖每秒。\omterm{def:g10:cell-metabolism:fermentation}{发酵}：$60/2 = 30$
分子葡萄糖每秒，是呼吸作用的十五倍。
\end{solution}

\begin{solution}{exo:g12:respiration-fermentation:10}
氰化物阻断呼吸链，电子无法传递给氧，
每分子葡萄糖约 30 个 \omterm{def:g12:respiration-fermentation:atp}{ATP} 中的 26 个随即丧失；\omterm{def:g10:cell-metabolism:fermentation}{发酵}
无法弥补这一缺口。脑和心脏消耗 \omterm{def:g12:respiration-fermentation:atp}{ATP} 最快，
而储备最少，因此数分钟内即会衰竭。
\end{solution}

\begin{solution}{exo:g12:respiration-fermentation:11}
二氧化碳释放于\omterm{def:g10:cells-common-unit:organelle}{线粒体}基质中，由\omterm{prop:g12:respiration-fermentation:mitochondrion}{克雷布斯循环}
（以及丙酮酸进入循环的步骤）产生；\omterm{prop:g12:respiration-fermentation:glycolysis}{糖酵解}中不产生二氧化碳。
水在呼吸链末端生成，
即电子和质子与氧结合之时。
\end{solution}

\begin{solution}{exo:g12:respiration-fermentation:12}
肌肉每\omterm{def:g10:cell-metabolism:fermentation}{发酵}一分子葡萄糖获得 2 个 \omterm{def:g12:respiration-fermentation:atp}{ATP}；由乳酸重新合成
该葡萄糖需肝脏消耗约 6 个 \omterm{def:g12:respiration-fermentation:atp}{ATP}，这些能量由其他燃料经
呼吸作用提供。机体为获得 2 个 \omterm{def:g12:respiration-fermentation:atp}{ATP} 花费了 6 个：\omterm{def:g10:cell-metabolism:fermentation}{发酵}是借用能量，
而借来的能量须连本带利偿还。
\end{solution}

\begin{solution}{exo:g12:respiration-fermentation:13}
在空气中，呼吸作用每分子葡萄糖产生 30 个 \omterm{def:g12:respiration-fermentation:atp}{ATP}，所需葡萄糖很少，
无须分流去生成乙醇；密封时，\omterm{def:g10:cell-metabolism:fermentation}{发酵}只产生 2 个，
因此消耗的葡萄糖多十五倍，并生成乙醇。面团
依靠二氧化碳膨胀，两条途径均会产生二氧化碳（呼吸作用每分子葡萄糖
产生 6 分子，\omterm{def:g10:cell-metabolism:fermentation}{发酵}产生 2 分子）；在致密的面团中，
氧气本就很快耗尽。
\end{solution}

\begin{solution}{exo:g12:respiration-fermentation:14}
热量：呼吸链的能量不再被捕获为 \omterm{def:g12:respiration-fermentation:atp}{ATP}，
而是直接释放。新生儿散热快，且不能很好地颤抖产热；棕色脂肪
是体内的“天然暖炉”。
\end{solution}

\begin{solution}{exo:g12:respiration-fermentation:15}
电子：\omterm{def:g10:cell-metabolism:autotrophy}{光合作用}从水中夺取电子并将其存入
糖中；呼吸作用从糖中取出电子并交给氧。
气体：\omterm{def:g10:cell-metabolism:autotrophy}{光合作用}消耗 $\mathrm{CO_2}$，释放
$\mathrm{O_2}$；呼吸作用则相反。能量：\omterm{def:g10:cell-metabolism:autotrophy}{光合作用}吸收
光能并储存于糖中；呼吸作用则从糖中释放能量，形成
\omterm{def:g12:respiration-fermentation:atp}{ATP} 和热量。\omterm{def:g10:cells-common-unit:organelle}{细胞器}：\omterm{def:g12:photosynthesis:chloroplast}{叶绿体}和\omterm{def:g10:cells-common-unit:organelle}{线粒体}。二者各自消耗
对方所产生的物质：\omterm{def:g10:cell-metabolism:autotrophy}{光合作用}会耗尽
$\mathrm{CO_2}$，呼吸作用会耗尽 $\mathrm{O_2}$ 和食物。
\end{solution}

\begin{solution}{pb:g12:respiration-fermentation:1}
\textbf{1.} 三分之一的 \qty{1000}{W} 为 \qty{300}{W}：每分钟 \qty{18}{kJ}，
即每分钟 \qty{0.6}{mol} \omterm{def:g12:respiration-fermentation:atp}{ATP}。

\textbf{2.} 每分钟约 \qty{300}{g} ——相比之下，体内仅储存几克：
储备被不断循环利用。

\textbf{3.} \qty{50}{g} 为 \qty{0.1}{mol}：每个 \omterm{def:g12:respiration-fermentation:atp}{ATP} 分子
每分钟约被再生 6 次。

\textbf{4.} 每分钟葡萄糖约为 $0.6/30 = \qty{0.02}{mol}$（约
\qty{3.6}{g}）。

\textbf{5.} 由 $6 \times 0.02 = \qty{0.12}{mol}$ 得 $\mathrm{O_2}$，约
每分钟 \qty{2.9}{L}，低于她的\omterm{def:g10:exercise-and-energy:vo2max}{最大摄氧量} \qty{4}{L/min}，且
与每升氧气 \qty{20}{kJ} 的数值相符，见
\cref{ch:g10:exercise-and-energy}：\qty{1000}{W} 的功率
可由呼吸作用维持。

\textbf{6.} \omterm{def:g10:cells-common-unit:cell}{细胞质}中 2 个，\omterm{def:g10:cells-common-unit:organelle}{线粒体}中 28 个：约占
93\% 的 \omterm{def:g12:respiration-fermentation:atp}{ATP}。

\textbf{7.} 由 $6 \times 0.02 = \qty{0.12}{mol}$ 得每分钟产生的 $\mathrm{CO_2}$
全部来自基质中的\omterm{prop:g12:respiration-fermentation:mitochondrion}{克雷布斯循环}与丙酮酸的进入。

\textbf{8.} $30 \times 30 = \qty{900}{kJ}$，约占 \qty{2870}{kJ} 的
三分之一；其余为热量，她通过出汗和
流向皮肤的血流将其散发。

\textbf{9.} 呼吸链的\omterm{def:g11:gene-expression:protein}{蛋白质}和合成 \omterm{def:g12:respiration-fermentation:atp}{ATP} 的酶位于
膜上；\omterm{def:g12:respiration-fermentation:atp}{ATP} 合成的速率与承载它们的
膜面积成正比。

\textbf{10.} 这提高了呼吸作用合成 \omterm{def:g12:respiration-fermentation:atp}{ATP} 的速率
（以及她能够燃烧的脂肪比例），从而提高她在不依赖\omterm{def:g10:cell-metabolism:fermentation}{发酵}的情况下
所能维持的功率；它不改变每分子葡萄糖的产量，也不改变
每分子葡萄糖所需的氧气量。

\textbf{11.} $800 \times 30 = \qty{24}{kJ}$：即 \qty{0.8}{mol} \omterm{def:g12:respiration-fermentation:atp}{ATP}。

\textbf{12.} 呼吸作用：$400 \times 30 = \qty{12}{kJ}$，即
\qty{0.4}{mol}；其余 \qty{0.4}{mol} 必须由\omterm{def:g10:cell-metabolism:fermentation}{发酵}提供。

\textbf{13.} 需要 $0.4/2 = \qty{0.2}{mol}$ 的葡萄糖，并产生
\qty{0.4}{mol} 乳酸。

\textbf{14.} 呼吸作用：每个 \omterm{def:g12:respiration-fermentation:atp}{ATP} 需 $1/30$ 个葡萄糖；\omterm{def:g10:cell-metabolism:fermentation}{发酵}：$1/2$
——葡萄糖多出十五倍。肌肉之所以接受，是因为此刻需要
\omterm{def:g12:respiration-fermentation:atp}{ATP}，而\omterm{def:g10:exercise-and-energy:glycogen}{糖原}就在那里；氧债稍后偿还。

\textbf{15.} 乳酸必须被氧化，或重新合成为葡萄糖，
而这一过程需消耗由呼吸作用产生的 \omterm{def:g12:respiration-fermentation:atp}{ATP}：耗氧量将持续高于
静息水平，直至氧债还清。

\textbf{16.} \qty{500}{W} 持续一小时为 \qty{1800}{kJ}：$1800/38
\approx \qty{47}{g}$ 的脂肪。

\textbf{17.} 脂肪的还原程度更高——富含氢——因此每
克脂肪完全氧化所需的氧气更多；每升氧气产生的能量
略少于葡萄糖。

\textbf{18.} \omterm{prop:g12:respiration-fermentation:glycolysis}{糖酵解}及其\omterm{def:g10:cell-metabolism:fermentation}{发酵}：脂肪不能\omterm{def:g10:cell-metabolism:fermentation}{发酵}，
因此依靠\omterm{def:g10:cell-metabolism:fermentation}{发酵}供能的短跑无法利用脂肪——只能利用
葡萄糖和\omterm{def:g10:exercise-and-energy:glycogen}{糖原}。

\textbf{19.} 肝脏释放自身\omterm{def:g10:exercise-and-energy:glycogen}{糖原}分解出的葡萄糖，并
利用乳酸重新合成部分葡萄糖，使血糖维持在 \qty{1}{g/L} 左右
（见\cref{ch:g12:glucose-and-diabetes}）；大脑的供应
得到保护。

\textbf{20.} 每分钟约 \qty{0.6}{mol} \omterm{def:g12:respiration-fermentation:atp}{ATP}；
每分钟约 \qty{3}{L} 氧气即可支撑这些 \omterm{def:g12:respiration-fermentation:atp}{ATP} 的生成；\omterm{def:g10:cells-common-unit:organelle}{线粒体}内膜
上的呼吸链承担其中约十分之九。
\end{solution}
